\documentclass[a4paper, 11pt]{article}

% -- Coments
\usepackage{verbatim}

% ----- Fonts -----
% -- Fuente --
% \usepackage{fontspec}
% \setmonofont{JetBrainsMono Nerd Font}

\usepackage{babel}
\babelprovide[import, main]{spanish}

% -- Color --
\usepackage{xcolor}
%\definecolor{azul}{RGB}{00,33,99}
\definecolor{azul}{RGB}{35,72,180}

% -- Page Margin --
\usepackage[margin=1in]{geometry}

% -- Espaciados --
\newcommand{\Pspace}{0.5cm}
\newcommand{\Aspace}{0.2cm}

% -- Columnas --
\usepackage{multicol}

% -- Imagenes --
\usepackage{graphicx}
\usepackage{float}

% -- Matemáticas --
\usepackage{amsmath, amssymb}

% -- Gráficas --
\usepackage{pgfplots}
\pgfplotsset{compat=1.18}


\title
{
    Contexto Regional 2026-2 \\
    El interés expansionista norteamericano en Sonora: 1848-1861 - Ana Rosa Suárez Argüello \\

    Profesor José Moreno Vega
}

\begin{document}

\begin{otherlanguage}{spanish}
\maketitle
\end{otherlanguage}

\begin{center}
    \begin{tabular}{r|l}
        \textbf{Expediente} & \textbf{Nombre} \\ \hline
        219208106 & Bórquez Guerrero Angel Fernando \\
    \end{tabular}
\end{center}

{\Large 1. ¿Cuál es el argumento central del autor?} \\
Habla sobre el interés de los Estados Unidos de América por quedarse con Sonora entre 1848 y 1861. No era puro afán de invadir, sino que había razones económicas de por medio, como los recursos agrícolas y mineros, además de rutas estratégicas como el ferrocarril y el puerto de Guaymas. La autora sostiene que esto no fue cosa de filibusteros sueltos, sino un sentimiento compartido por la gente, por los californianos y por el mismo gobierno, impulsado por el "destino manifiesto" y la fama de la riqueza sonorense. Los filibusteros (Morehead, Walker, Crabb) fueron solo la parte más visible de ese ánimo, y el gobierno lo apoyó con presión diplomática, amenazas de usar la fuerza y chantaje económico.\par
\vspace{1cm}

{\Large 2. ¿Qué evidencias utiliza para respaldarlo?} \par
\begin{itemize}
    \item Prensa y opiniones públicas.
    \item Documentos y acciones del gobierno estadounidense.
    \item El tratado de La Mesilla donde se intentó comprar más territorio por parte del gobierno de los Estados Unidos.
\end{itemize}


\end{document}
